\section{Exact two-level factorization}
\label{sec:finite-factorization}

Let \(a=(a_n)_{n\in\Z}\) be finitely supported and complex.  Define
\begin{equation}
 S=\#\supp a,\qquad
 L=\|a\|_\infty,\qquad
 M=\|a\|_1,\qquad
 D=\|a\|_2^2,\qquad
 Z=\sum_na_n.
\label{eq:SLMDZ}
\end{equation}

\subsection{The zero sequence}

\begin{lemma}[Edge equivalence]
\label{lem:zero-edge}
The following are equivalent:
\[
 a=0,\qquad M=0,\qquad D=0.
\]
Whenever they hold, \(S=L=Z=0\).
\end{lemma}

\begin{proof}
Both \(M\) and \(D\) are sums of nonnegative terms and vanish exactly
when every coefficient vanishes.  The remaining conclusions are
then immediate.
\end{proof}

For \(a=0\), ratios such as \(M^2/D\) and \(|Z|/M\) are formally
\(0/0\).  We remove the ambiguity by convention:
\begin{equation}
 \canc(0)=0,\qquad
 \Seff(0)=0,\qquad
 \Flat(0)=0.
\label{eq:zero-convention}
\end{equation}
This agrees with the zero-flatness convention used in the physical
route.  All division in the remainder of the finite theory is
restricted to \(a\ne0\).

\subsection{The factorization}

\begin{definition}[Two levels]
\label{def:two-level-statistics}
For \(a\ne0\), define
\begin{equation}
 \canc(a)=\frac{|Z|}{M},\qquad
 \Seff(a)=\frac{M^2}{D},\qquad
 \Flat(a)=\frac{|Z|^2}{D}.
\label{eq:two-level-statistics}
\end{equation}
We call \(\canc\) the relative phase-cancellation coefficient,
\(\Seff\) the effective support, and \(\Flat\) the zero-frequency
concentration.
\end{definition}

\begin{theorem}[Exact mass--cancellation factorization]
\label{thm:exact-factorization}
For every finitely supported complex sequence,
\begin{equation}
 \boxed{\Flat(a)=\canc(a)^2\Seff(a).}
\label{eq:exact-factorization}
\end{equation}
For \(a\ne0\), the elementary ranges are
\begin{equation}
 0\le\canc(a)\le1,\qquad
 1\le\Seff(a)\le S,\qquad
 0\le\Flat(a)\le\Seff(a)\le S.
\label{eq:elementary-ranges}
\end{equation}
All three statistics are invariant under multiplication of \(a\)
by a nonzero complex scalar.
\end{theorem}

\begin{proof}
For \(a=0\), \eqref{eq:exact-factorization} is
\(0=0^2\cdot0\) by convention.  Suppose \(a\ne0\).
Then \(M,D>0\), and direct multiplication gives
\[
 \canc^2\Seff
 =
 \frac{|Z|^2}{M^2}\frac{M^2}{D}
 =
 \frac{|Z|^2}{D}
 =
 \Flat.
\]
The triangle inequality gives \(|Z|\le M\), hence
\(0\le\canc\le1\).  Also
\[
 D\le M^2
\]
and Cauchy--Schwarz on the exact support gives
\[
 M^2\le SD.
\]
Thus \(1\le\Seff\le S\).  The remaining bounds follow from the
factorization.  Each numerator and denominator has the same
homogeneity needed for scalar invariance.
\end{proof}

\begin{remark}[Magnitude level versus phase level]
\label{rem:magnitude-phase}
The statistic \(\Seff\) depends only on
\((|a_n|)\).  For fixed magnitudes, changing phases leaves
\(\Seff\), \(M\), and \(D\) unchanged but can change \(\canc\) and
\(\Flat\).  This is the exact sense in which the factorization has
two levels.
\end{remark}

\subsection{Equality cases}

\begin{theorem}[Equality dictionary]
\label{thm:equality-dictionary}
Let \(a\ne0\).
\begin{enumerate}[label=(\roman*),leftmargin=2.2em]
\item \(\canc=1\) if and only if all nonzero coefficients lie on one
closed complex ray: there is \(\zeta\in\C\), \(|\zeta|=1\), such
that
\[
 a_n=\zeta|a_n|
 \qquad(n\in\supp a).
\]
\par
\item \(\canc=0\) and \(\Flat=0\) if and only if \(Z=0\).
\item \(\Seff=S\) if and only if \(|a_n|\) is constant on
\(\supp a\).
\item \(\Seff=1\) if and only if \(S=1\).
\item \(\Flat=\Seff\) if and only if the common-ray condition in
(i) holds.
\item \(\Flat=S\) if and only if \(a_n\) is one fixed nonzero complex
number on \(\supp a\).
\end{enumerate}
\end{theorem}

\begin{proof}
Part (i) is the equality case of the complex triangle inequality.
Part (ii) follows from \(D,M>0\).  Part (iii) is the equality case of
\[
 \left(\sum_{n\in\supp a}|a_n|\right)^2
 \le S\sum_{n\in\supp a}|a_n|^2.
\]
For (iv), expand
\[
 M^2
 =D+2\sum_{m<n}|a_m||a_n|.
\]
Equality \(M^2=D\) holds exactly when at most one coefficient is
nonzero.  Parts (v) and (vi) follow from
\(\Flat=\canc^2\Seff\): equality in (v) requires
\(\canc=1\), while equality \(\Flat=S\) requires both
\(\canc=1\) and \(\Seff=S\).  The latter two conditions say that
all moduli and all phases agree.
\end{proof}
